\chapter{VC Dimension and Uniform Learning}
\label{sec:p15-vc-dimension-and-uniform-learning}

Pointwise concentration for a fixed hypothesis does not justify selecting
that hypothesis after inspecting the data. Uniform control asks for one
event on which every candidate's empirical frequency is close to its
probability. The classical question is identified by
\cite{VapnikChervonenkis1971Uniform}. The capacity/ghost-sample construction
below is a fully proved editorial route, with an explicit measurability
hypothesis. Its relationship to the primary author presentation
\cite[Sections4--5.6]{LuxburgScholkopf2008Author} is stated without equating
different source versions, printed constants or consistency quantifiers.

\section{Traces, shattering and a measurable supremum}
\label{sec:p15-vcl-capacity}

\begin{hbdefinition}{def:p15-vcl-event-class}{Event class, trace growth and discrepancy}
Let \((Z,\mathcal A,P)\) be a probability space and let
\(\mathcal C\subseteq\mathcal A\) be a nonempty class fixed before
sampling. Its trace on a finite set \(S\subseteq Z\) is
\(\mathcal C\upharpoonright S=\{C\cap S:C\in\mathcal C\}\). It shatters \(S\)
when that trace is \(\mathcal P(S)\). Its VC dimension is the supremum
of the sizes of shattered finite sets, possibly infinity; the empty set
is shattered, so dimension zero is meaningful.

For \(t\geq0\), let \(\Pi_{\mathcal C}(t)\) be the largest number
of distinct vectors \((\mathbf1_C(z_1),\ldots,\mathbf1_C(z_t))\)
on any ordered \(t\)-tuple. Repetitions are allowed and contribute no
independent coordinate choices. \(\Pi_{\mathcal C}(0)=1\).
For an iid sample \(Z_1,\ldots,Z_m\) of law \(P\), \(m\geq1\),
write
\[
 P_m(C)=m^{-1}\sum_i\mathbf1_C(Z_i),\qquad
 \Delta_m=\sup_{C\in\mathcal C}|P_m(C)-P(C)|.
\]
A supremum over an uncountable class is not automatically measurable.
Absolute bars around a scalar deviation mean numerical absolute value;
cardinality bars around a finite set still count its members.
The author presentation's \(N(F,n)\) is the corresponding trace-growth
quantity; its labels \(-1,+1\) correspond bijectively to \(0,1\) here.
\end{hbdefinition}

\begin{hbdefinition}{def:p15-vcl-countable-core}{Pointwise measurable core}
Assume a nonempty countable \(\mathcal C_0\subseteq\mathcal C\)
such that each \(\mathbf1_C\) is a pointwise limit of a sequence
of indicators from \(\mathcal C_0\). This is a genuine approximation
condition on individual functions, not an assumption of uniform convergence
or its probability bound. Countable classes satisfy it trivially.
\end{hbdefinition}

Finite VC dimension does not itself supply a pointwise countable core.
All singleton sets and the empty set on an uncountable domain have dimension
one, but a countable subfamily has only countably many points in its union.
At a point outside that union its indicators cannot converge to that point's
singleton indicator. Such a class needs a different established
measurability/admissibility or outer-probability treatment; no ordinary
probability assertion for its uncountable suprema is presumed here.

\begin{hblemma}{thm:p15-vcl-core}{proved=annex:p15.vcl.core}{Measurability and finite traces from the core}
Under the core condition, the suprema defining \(\Delta_m\) and the
two-sample discrepancy equal those over \(\mathcal C_0\), hence are
measurable. Every finite trace of \(\mathcal C\) occurs already in
\(\mathcal C_0\). The event that some \(C\) has \(P_m(C)=0\)
and \(P(C)>\varepsilon\) is also the corresponding countable union.
\end{hblemma}

\begin{hbtheorem}{thm:p15-vcl-sauer}{proved=annex:p15.vcl.sauer}{The full Sauer trace bound}
If the VC dimension is \(d<\infty\), then for every \(t\geq0\),
\[
 \Pi_{\mathcal C}(t)\leq\sum_{j=0}^{\min(d,t)}\binom tj.
\]
For \(1\leq d\leq t\), this is at most \((et/d)^d\).
For \(d=0\), every indicator in the nonempty class is the same
function and \(\Pi_{\mathcal C}(t)=1\). If \(d=\infty\),
\(\Pi_{\mathcal C}(t)=2^t\) for every \(t\).
\end{hbtheorem}

This counts traces on sampled points, not all hypotheses in the class.
It applies to infinite and, under the core condition, uncountable classes.
No finite-cardinality union bound on the original class is substituted for it.

\section{The ghost sample and randomized split}
\label{sec:p15-vcl-uniform}

\begin{hbtheorem}{thm:p15-vcl-uniform}{proved=annex:p15.vcl.uniform}{Uniform discrepancy with explicit constants}
For the fixed event class with a countable core, iid sampling gives, for
\(m\geq1\) and \(0<\varepsilon<1\),
\[
 P^{\otimes m}(\Delta_m>\varepsilon)
 \leq\min\{1,\ 4\Pi_{\mathcal C}(2m)e^{-m\varepsilon^2/8}\}.
\]
The nontrivial ghost step uses \(m\varepsilon^2\geq2\); below that
threshold the displayed minimum is justified by the trivial bound one.
In dimension zero, the stronger pointwise bound
\(2e^{-2m\varepsilon^2}\) applies. Finite VC dimension therefore gives
uniform convergence in probability, uniformly over sampling laws \(P\)
for which the stated measurability hypothesis holds.
\end{hbtheorem}

The annex selects a measurable bad witness, introduces an independent
sample, fixes the pooled observations and randomly swaps each pair before
using concentration. Only then is the finite trace union taken. Treating
a data-dependent trace representative as though it were a fixed hypothesis
under the original sampling law would omit that step. The constants above
belong to this proof; they are not silently identified with the different
printed constants in the 2008 author's explanatory calculation.

\section{A measurable learner and the two risk regimes}
\label{sec:p15-vcl-learning}

\begin{hbdefinition}{def:p15-vcl-classification}{Binary loss and realizability}
Let \(P\) now be a joint law on \(X\times\{0,1\}\), and let
\(H\) be a nonempty predeclared class of measurable binary predictors
with a pointwise countable core \(H_0=(h_j)\). Set
the positive prediction-event class to \(\{\{x:h(x)=1\}:h\in H\}\);
\(\Pi_H\) and its dimension mean the growth and dimension of this class.
The loss-event class is
\(\mathcal E_H=\{\{(x,y):h(x)\ne y\}:h\in H\}\). Define
\[
 R(h)=P\{h(X)\ne Y\},\qquad
 \widehat R_m(h)=m^{-1}\sum_i\mathbf1_{\{h(X_i)\ne Y_i\}},\qquad
 R_H^*=\inf_{h\in H}R(h).
\]
The regime is realizable if some \(h\in H\) has risk zero. Agnostic
learning makes no such assumption; labels can be noisy and the class can
miss the best predictor. Neither \(R_H^*=0\) without attainment nor a
small observed error is silently called realizability.
\end{hbdefinition}

\begin{hbtheorem}{thm:p15-vcl-erm}{proved=annex:p15.vcl.erm}{An actual measurable empirical minimizer}
The binary error-event class has the same trace growth as \(H\).
Choose the first core index attaining the minimum empirical error.
That index exists on every sample and is measurable; it gives an exact
ERM over \(H\), though no terminating effective search is promised.
An \(\eta\)-approximate ERM, \(\eta\geq0\), means
\(\widehat R_m(\widehat h)\leq\inf_{h\in H}\widehat R_m(h)+\eta\).
On the event of uniform risk discrepancy at most \(u\), every
\(\eta\)-approximate ERM satisfies
\[
 R(\widehat h)\leq R_H^*+2u+\eta.
\]
In the realizable regime an exact ERM is consistent with the sample almost
surely. These are excess-risk conclusions within \(H\), not automatic
Bayes-risk or true-target identification guarantees.

If a training-derived predictor is jointly measurable and an independent
validation sample of size \(v\geq1\) has the same law \(P\), then
\(P(|\widehat R_v(\widehat h)-R(\widehat h)|>u)\leq2e^{-2vu^2}\).
The predictor must not be selected further using that validation sample
under this one-predictor statement.
\end{hbtheorem}

\begin{hbtheorem}{thm:p15-vcl-consistent}{proved=annex:p15.vcl.consistent}{The faster consistent-learning sampling argument}
For the same loss class, if \(m\varepsilon\geq8\),
\[
 P\{\exists h\in H:\widehat R_m(h)=0, R(h)>\varepsilon\}
 \leq2\Pi_H(2m)2^{-m\varepsilon/2}.
\]
Thus every measurable consistent learner has this error guarantee in a
realizable problem. It does not guarantee that a consistent learner exists
under arbitrary noise. In dimension zero and a realizable problem, the
sole prediction function has risk zero, without this sample-size restriction.
\end{hbtheorem}

\begin{hbcorollary}{thm:p15-vcl-samples}{proved=annex:p15.vcl.samples}{Finite sufficient sample sizes, with no implicit log-sample equation}
For \(d\geq1\), \(0<\varepsilon,\delta<1\), a sufficient sample size
for uniform discrepancy at most \(\varepsilon\) with failure probability
at most \(\delta\) is
\[
 m\geq\left\lceil\frac{64}{\varepsilon^2}
       \left[d\ln\frac{64}{\varepsilon^2}+\ln\frac8\delta\right]\right\rceil.
\]
For consistent realizable learning, a sufficient size for risk at most
\(\varepsilon\) is
\[
 m\geq\left\lceil\frac{32}{\varepsilon}
       \left[d\ln\frac{64}{\varepsilon}+\ln\frac4\delta\right]\right\rceil.
\]
For exact agnostic ERM and desired excess risk \(\tau\in(0,1)\), use
the first formula with \(\varepsilon=\tau/2\). The constants are
sufficient, not sharp or necessary. No empirical mean or denominator is
defined here for an empty sample.
\end{hbcorollary}

\section{Worked infinite classes and failure boundaries}
\label{sec:p15-vcl-examples}

\begin{hbtheorem}{thm:p15-vcl-examples}{proved=annex:p15.vcl.examples}{Capacity, measurability and noise are different conditions}
\begin{enumerate}
\item Real thresholds \(h_t(x)=\mathbf1_{\{x\leq t\}}\), together
with the two constant functions, have a rational-threshold countable core,
VC dimension one and growth \(t+1\) on \(t\) distinct real points.
\item All finite subsets of \(\mathbb N\) form a countable event class
of infinite VC dimension. For each sample size \(m\), a uniform law on
\(2m\) points has discrepancy at least \(1/2\) on every sample.
For a point-mass law the same class has discrepancy zero almost surely.
Thus a distribution-free failure is not a fixed-law necessity claim.
\item For uniform \(X\in[0,1]\) and \(Y=\mathbf1_{\{X\geq1/2\}}\),
the classifier that memorizes the training labels and predicts one elsewhere
is jointly measurable, has training error zero and true risk \(1/2\).
This is the actual memorization case in the 2008 author presentation,
with its coding of labels translated explicitly to \(0,1\).
\item If \(Y\) is a fair bit independent of \(X\), every deterministic
predictor has risk \(1/2\). The threshold class is still finite-dimensional
and agnostic ERM still controls excess risk, but the problem is not realizable.
\end{enumerate}
\end{hbtheorem}

Independent sampling, a fixed class and measurable selection remain real
hypotheses. The repeated-input and distribution-shift countercases in
Theorem~\ref{thm:p15-occam-boundaries} do not acquire an iid guarantee merely
because the class has low VC dimension. A validation claim concerns its
specified future law, not a changed deployment law.
The finite-class and code-weighted results in
Theorems~\ref{thm:p15-conc-finite-learning} and~\ref{thm:p15-occam-compression}
have different complexity quantities and realizability hypotheses.
Here the class may be infinite; it is its sampled traces that become finite.

This construction proves a measurable capacity-based sufficient route and
its learning consequences. It does not assert an unqualified converse for
every fixed distribution or every ERM tie policy, an arbitrary real-loss
analogue, or the full source-specific entropy/necessary-and-sufficient
classification of the selected 1971 presentation.
